\documentclass[final]{beamer}
\usepackage[orientation=portrait,size=a0,scale=1.4]{beamerposter}
\usepackage{graphicx}
\usepackage{booktabs}
\usepackage{tikz}
\usepackage{pgfplots}
\pgfplotsset{compat=1.16}
\usetikzlibrary{positioning,arrows.meta,calc}
\usepackage{enumitem}

\usetheme{default}
\usecolortheme{whale}
\setbeamercolor{block title}{fg=white,bg=blue!70!black}
\setbeamercolor{block body}{fg=black,bg=blue!5}
\setbeamercolor{block alerted title}{fg=white,bg=red!70!black}
\setbeamercolor{block alerted body}{fg=black,bg=red!5}
\setbeamerfont{block title}{size=\large,series=\bfseries}
\setbeamertemplate{navigation symbols}{}

\title{\Huge Cut Rules in DAG-Based BFT Mempools:\\From Relaxed Mempool to Adaptive Cuts}
\author{\Large Woochang Jeong \and Chanik Park}
\institute{\large Computer Science and Engineering, POSTECH, South Korea}
\date{}

\begin{document}
\begin{frame}[t]

\vspace{-1cm}
\begin{beamercolorbox}[wd=\paperwidth,colsep*=1.5cm]{title}
\usebeamerfont{title}\usebeamercolor[fg]{title}\inserttitle\par
\vspace{0.5cm}
\usebeamerfont{author}\usebeamercolor[fg]{author}\insertauthor\par
\vspace{0.3cm}
\usebeamerfont{institute}\usebeamercolor[fg]{institute}\insertinstitute
\vspace{0.5cm}
\end{beamercolorbox}

\vspace{1cm}

\begin{columns}[T]

%=== LEFT COLUMN ===
\begin{column}{0.48\textwidth}

\begin{block}{1. Motivation: The Cut Rule Problem}
\large
DAG-based BFT protocols decouple data dissemination from ordering. Validators broadcast batches into a DAG; a consensus leader periodically \textbf{cuts} the DAG to commit transactions.

\vspace{0.5cm}
\textbf{Key Question:} How should the leader decide \emph{which} data to include in each commit?

\vspace{0.5cm}
The \textbf{cut rule} determines:
\begin{itemize}[leftmargin=2cm]
    \item[\textbullet] \textbf{Straggler resilience:} Do slow validators delay commits?
    \item[\textbullet] \textbf{Transaction fairness:} Are all validators' data included?
    \item[\textbullet] \textbf{Hangover effect:} How fast does the system recover after blips?
\end{itemize}
\end{block}

\vspace{1cm}

\begin{block}{2. Two Existing Cut Strategies}
\large

\vspace{0.5cm}
\begin{center}
\begin{tikzpicture}[
    scale=1.8, transform shape,
    block/.style={rectangle, draw, minimum width=1.2cm, minimum height=0.8cm, font=\small},
    committed/.style={block, fill=green!30},
    uncommitted/.style={block, fill=red!20},
    label/.style={font=\Large\bfseries},
]
% Left: Conservative cut
\node[label] at (3, 5.5) {Conservative Cut};
\node[label, blue!70!black] at (3, 4.8) {(Relaxed Mempool)};

\node[font=\large] at (-0.8, 4) {$v_1$};
\node[font=\large] at (-0.8, 3) {$v_2$};
\node[font=\large] at (-0.8, 2) {$v_3$};
\node[font=\large] at (-0.8, 1) {$v_4$};

\foreach \x in {1,...,5} { \node[committed] at (\x*1.1, 4) {}; }
\foreach \x in {1,...,5} { \node[committed] at (\x*1.1, 3) {}; }
\foreach \x in {1,...,3} { \node[committed] at (\x*1.1, 2) {}; }
\node[uncommitted] at (4.4, 2) {};
\foreach \x in {1,...,2} { \node[uncommitted] at (\x*1.1, 1) {}; }

\draw[dashed, ultra thick, blue] (3.7, 0.4) -- (3.7, 4.6);
\node[font=\Large, blue] at (3.7, 0.0) {$H_{\min}=3$};

% Right: Inclusive cut
\node[label] at (11, 5.5) {Inclusive Cut};
\node[label, orange!70!black] at (11, 4.8) {(Autobahn)};

\node[font=\large] at (7.2, 4) {$v_1$};
\node[font=\large] at (7.2, 3) {$v_2$};
\node[font=\large] at (7.2, 2) {$v_3$};
\node[font=\large] at (7.2, 1) {$v_4$};

\foreach \x in {1,...,5} { \node[committed] at (7.2+\x*1.1, 4) {}; }
\foreach \x in {1,...,5} { \node[committed] at (7.2+\x*1.1, 3) {}; }
\foreach \x in {1,...,4} { \node[committed] at (7.2+\x*1.1, 2) {}; }
\foreach \x in {1,...,2} { \node[committed] at (7.2+\x*1.1, 1) {}; }

\draw[->, ultra thick, orange] (12.9, 3.5) -- (12.9, 4.4);
\draw[->, ultra thick, orange] (12.9, 2.5) -- (12.9, 3.4);
\draw[->, ultra thick, orange] (11.8, 1.5) -- (11.8, 2.4);
\draw[->, ultra thick, orange] (9.5, 0.5) -- (9.5, 1.4);
\end{tikzpicture}
\end{center}

\vspace{0.5cm}
\begin{center}
\begin{tabular}{@{}p{0.42\textwidth}p{0.42\textwidth}@{}}
\textbf{\color{blue!70!black}Conservative (Relaxed Mempool):} &
\textbf{\color{orange!70!black}Inclusive (Autobahn):} \\[0.3cm]
Top $2f{+}1$ validators by height &
All $n$ lanes included \\
Uniform $H_{\min}$ cut &
Per-lane tip references \\
Slow nodes \textbf{excluded} from cut &
\textbf{No node excluded} \\
Data sync before finalization &
Data sync after voting \\
\textcolor{green!50!black}{+ Low straggler latency} &
\textcolor{green!50!black}{+ Fast blip recovery} \\
\textcolor{red!70!black}{-- Slow blip recovery} &
\textcolor{red!70!black}{-- Higher straggler latency} \\
\end{tabular}
\end{center}
\end{block}

\vspace{1cm}

\begin{block}{3. Our Proposal: Adaptive Cut}
\large

\vspace{0.3cm}
\textbf{Key Insight:} The two strategies have \emph{complementary} strengths. We unify them with an \textbf{adaptive cut} that monitors \emph{height disparity} ($\Delta = h_{\max} - h_{\min}$) across validators.

\vspace{0.5cm}
\begin{center}
\begin{tikzpicture}[scale=1.5, transform shape,
    box/.style={rectangle, draw, rounded corners, minimum width=4cm, minimum height=1.2cm, font=\large, thick},
]
\node[box, fill=green!20] (monitor) at (0, 0) {Monitor $\Delta = h_{\max} - h_{\min}$};
\node[box, fill=blue!15] (relaxed) at (-4.5, -2.5) {Conservative Cut};
\node[box, fill=orange!15] (inclusive) at (4.5, -2.5) {Inclusive Cut};

\draw[->, ultra thick] (monitor) -- node[left, font=\normalsize, align=right] {$\Delta > \delta$\\(stragglers)} (relaxed);
\draw[->, ultra thick] (monitor) -- node[right, font=\normalsize, align=left] {$\Delta \leq \delta$\\(converging)} (inclusive);

\node[font=\normalsize, align=center] at (-4.5, -4) {Low latency\\Exclude slow nodes};
\node[font=\normalsize, align=center] at (4.5, -4) {Fast recovery\\Commit all backlog};
\end{tikzpicture}
\end{center}

\vspace{0.3cm}
Both strategies operate on the \textbf{same DAG substrate}---switching is zero-overhead.
\end{block}

\end{column}

%=== RIGHT COLUMN ===
\begin{column}{0.48\textwidth}

\begin{block}{4. Experimental Results}
\large
\textbf{Setup:} 4 servers (i9-13900K, 64\,GB), 10 validators, 512\,B tx, 500\,KB batches.

\vspace{0.5cm}
\textbf{(a) Straggler Scenario} (4 validators with +300\,ms delay)

\vspace{0.3cm}
\begin{center}
\begin{tikzpicture}[scale=2.2, transform shape]
\begin{axis}[
    xlabel={Throughput (Ktx/s)},
    ylabel={Latency (ms)},
    xmin=0, xmax=200,
    ymin=0, ymax=4500,
    legend pos=north west,
    legend style={font=\normalsize},
    grid=major,
    width=10cm, height=7cm,
]
\addplot[color=blue, mark=square*, thick, mark size=2pt] coordinates {
    (20, 800) (50, 1200) (80, 2500) (98, 4000)
};
\addplot[color=red, mark=triangle*, thick, mark size=2pt] coordinates {
    (20, 500) (50, 620) (100, 780) (150, 950) (170, 1200)
};
\addplot[color=orange, mark=o, thick, mark size=2pt] coordinates {
    (20, 520) (50, 680) (100, 850) (140, 1100) (160, 1500)
};
\addplot[color=green!60!black, mark=diamond*, thick, mark size=2pt] coordinates {
    (20, 500) (50, 610) (100, 770) (150, 940) (172, 1180)
};
\legend{Bullshark, Cons.\ (RM), Incl.\ (AB), \textbf{Adaptive}}
\end{axis}
\end{tikzpicture}
\end{center}

\vspace{0.2cm}
Adaptive $\approx$ Conservative under stragglers (high $\Delta$ triggers Relaxed mode).

\vspace{1cm}
\textbf{(b) Hangover Effect} (5s network blip at $t$=10s)

\vspace{0.3cm}
\begin{center}
\begin{tikzpicture}[scale=2.2, transform shape]
\begin{axis}[
    xlabel={Time (s)},
    ylabel={Latency (ms)},
    xmin=0, xmax=40,
    ymin=0, ymax=6000,
    legend pos=north east,
    legend style={font=\normalsize},
    grid=major,
    width=10cm, height=7cm,
]
\fill[gray!20] (axis cs:10,0) rectangle (axis cs:15,6000);
\addplot[color=blue, mark=none, ultra thick] coordinates {
    (0,600) (5,620) (10,650) (15,5500) (17,4200) (20,3000) (25,2000) (30,1200) (35,700) (40,620)
};
\addplot[color=red, mark=none, ultra thick, dashed] coordinates {
    (0,500) (5,520) (10,530) (15,4800) (17,3500) (20,2200) (25,1400) (30,800) (35,550) (40,510)
};
\addplot[color=orange, mark=none, ultra thick, dotted] coordinates {
    (0,510) (5,530) (10,540) (15,4500) (17,1800) (20,700) (25,550) (30,520) (35,510) (40,510)
};
\addplot[color=green!60!black, mark=none, ultra thick, dashdotted] coordinates {
    (0,500) (5,520) (10,530) (15,4600) (17,1900) (20,720) (25,540) (30,510) (35,510) (40,510)
};
\legend{Bullshark, Cons.\ (RM), Incl.\ (AB), \textbf{Adaptive}}
\node[font=\large, fill=white, inner sep=2pt] at (axis cs:12.5,5700) {\textbf{Blip}};
\end{axis}
\end{tikzpicture}
\end{center}

\vspace{0.2cm}
Adaptive $\approx$ Inclusive during recovery ($\Delta$ drops $\rightarrow$ switches to Inclusive mode).
\end{block}

\vspace{1cm}

\begin{alertblock}{5. Key Takeaways}
\large

\vspace{0.5cm}
\begin{center}
\begin{tabular}{@{}lccc@{}}
\toprule
\textbf{Scenario} & \textbf{Cons.\ (RM)} & \textbf{Incl.\ (AB)} & \textbf{Adaptive} \\
\midrule
Straggler (4 nodes) & \textcolor{green!50!black}{\textbf{170K tx/s}} & 160K tx/s & \textcolor{green!50!black}{\textbf{172K tx/s}} \\
Straggler latency   & \textcolor{green!50!black}{\textbf{1.2s}} & 1.5s & \textcolor{green!50!black}{\textbf{1.18s}} \\
Blip recovery time  & 20s & \textcolor{green!50!black}{\textbf{2s}} & \textcolor{green!50!black}{\textbf{2.5s}} \\
Happy path          & 195K & 205K & 200K \\
\bottomrule
\end{tabular}
\end{center}

\vspace{0.5cm}
\textbf{The cut rule is a first-class design dimension in DAG-based BFT.}\\
The adaptive cut achieves the best of both worlds by dynamically switching based on validator height disparity---\textbf{zero overhead}, same DAG substrate.

\vspace{0.5cm}
\end{alertblock}

\vspace{1cm}

\begin{block}{References}
\normalsize
{[1]} Danezis et al., ``Narwhal and Tusk,'' \emph{EuroSys}, 2022.\quad
{[2]} Spiegelman et al., ``Bullshark,'' \emph{ACM CCS}, 2022.\\
{[3]} Giridharan et al., ``Autobahn: Seamless high speed BFT,'' \emph{SOSP}, 2024.\quad
{[4]} Jeong \& Park, ``Relaxed Mempool for DAG-based BFT,'' 2025.
\end{block}

\end{column}
\end{columns}

\end{frame}
\end{document}
